\section{Introduction}
\label{sec:intro}

Automated analysis of laboratory animal behaviour has become an important
tool in biomedical research.  Reliably localising experimental mice in
diverse cage environments---with cluttered backgrounds, partial occlusion,
and varying lighting---is a prerequisite for downstream behaviour
quantification tasks such as trajectory tracking and posture estimation.

Classical frame-differencing or background-subtraction methods degrade in
the presence of illumination changes or when multiple animals interact.
Modern deep-learning object detectors offer superior robustness, but their
real-world deployment raises two practical questions that this project
addresses:

\begin{enumerate}[leftmargin=1.5em]
  \item \textbf{Training efficiency:} Given a limited GPU budget (dual
    Tesla T4, no InfiniBand), how do batch size, learning rate, number of
    GPUs, and dataset scale jointly affect final accuracy?

  \item \textbf{Mobile deployment:} Can the trained model run on-device
    (iPhone) at interactive frame rates without proprietary hardware
    accelerators?
\end{enumerate}

\subsection*{Contributions}

\begin{itemize}[leftmargin=1.5em]
  \item We curate and release a 7,286-image two-class mouse detection
    dataset assembled from three heterogeneous sources, providing train/val
    splits and VOC-format annotations.

  \item We design and execute a $4\!\times\!2$ controlled experiment matrix
    with learning rates derived from the Linear Scaling Rule, isolating the
    independent effects of GPU count, batch size, training duration, and
    dataset scale across eight training runs.

  \item We identify and quantify a \emph{counter-intuitive accuracy
    inversion}: small-batch single-GPU training (mAP = 94.30\%) surpasses
    the standard dual-GPU configuration (mAP = 93.80\%), and explain this
    through the implicit regularisation effect of gradient noise on
    small datasets.

  \item We document a complete deployment pipeline from PaddlePaddle
    training through ONNX export to React Native iOS inference, recording
    the four engineering iterations that raised throughput from 1\,FPS to
    14\,FPS.

  \item We demonstrate a striking data-scale effect on YOLOv3:
    a 3$\times$ increase in training images yields a 47-percentage-point
    jump in mAP, providing a concrete quantification of the
    ``data-first'' principle for anchor-based detectors.
\end{itemize}

\subsection*{Paper Organisation}

\Cref{sec:dataset} describes data collection and preprocessing.
\Cref{sec:methodology} reviews the two model architectures.
\Cref{sec:experiments} details the experimental design.
\Cref{sec:results} presents quantitative results and analysis.
\Cref{sec:deployment} covers the iOS deployment pipeline.
\Cref{sec:future} outlines future work.
\Cref{sec:conclusion} concludes.
